\documentclass[12pt,a4paper]{article}

\usepackage[margin=1in]{geometry}
\usepackage{amsmath}
\usepackage{booktabs}
\usepackage{siunitx}
\usepackage{enumitem}
\usepackage[T1]{fontenc}
\usepackage{lmodern}

\title{\textbf{Determination of Refractive Index and Molecular Refractivity}}
\author{}
\date{}

\begin{document}

\maketitle

\section*{Aim}

To determine the refractive indices of the given liquids: water, ethanol, glycerol, and cyclohexane using an Abbe refractometer, and to calculate their molecular refractivities.

\section*{Principle}

The refractive index of a liquid is the ratio of the velocity of light in vacuum to its velocity in the liquid:

\[
n = \frac{c}{v}
\]

An Abbe refractometer determines the refractive index from the critical angle at the interface between the sample and a high-refractive-index prism. The boundary between the illuminated and dark regions in the telescope is adjusted to the cross-wire, and the refractive index is read directly from the instrument.

The molecular refractivity is calculated using the Lorentz--Lorenz equation:

\[
R_m =
\left(\frac{n^2-1}{n^2+2}\right)
\frac{M}{\rho}
\]

where:
\begin{itemize}
    \item $R_m$ = molecular refractivity, in \si{\cubic\centi\meter\per\mole}
    \item $n$ = refractive index of the liquid
    \item $M$ = molar mass of the liquid, in \si{\gram\per\mole}
    \item $\rho$ = density of the liquid, in \si{\gram\per\cubic\centi\meter}
\end{itemize}

Thus, once $n$, $M$, and $\rho$ are known, molecular refractivity can be calculated.

\section*{Apparatus and Chemicals Required}

\subsection*{Apparatus}
\begin{itemize}
    \item Abbe refractometer
    \item Dropper
    \item Soft lens tissue or lint-free tissue
    \item Distilled water
\end{itemize}

\subsection*{Chemicals}
\begin{itemize}
    \item Water
    \item Ethanol
    \item Glycerol
    \item Cyclohexane
\end{itemize}

\section*{Procedure}

\begin{enumerate}
    \item Clean the prism surfaces of the Abbe refractometer carefully using a suitable solvent and soft lens tissue.
    \item Place a few drops of the liquid sample on the lower prism so that the entire prism surface is covered.
    \item Close the upper prism gently to form a thin, uniform layer of the liquid between the prisms.
    \item Direct the light source toward the refractometer.
    \item Look through the eyepiece and adjust the illumination so that a distinct light--dark boundary is visible.
    \item Adjust the dispersion compensator, if provided, to remove colour fringes and obtain a sharp boundary.
    \item Rotate the measuring control until the light--dark boundary coincides exactly with the cross-wire.
    \item Read the refractive index from the refractometer scale.
    \item Record the refractive index of the liquid.
    \item Clean the prisms thoroughly before measuring the next liquid.
    \item Repeat the procedure for water, ethanol, glycerol, and cyclohexane.
    \item Using the measured refractive index and the density and molar mass of each liquid, calculate molecular refractivity using the Lorentz--Lorenz equation.
\end{enumerate}

\section*{Observation Table}

\begin{table}[h]
\centering
\caption{Refractive indices and molecular refractivities of the given liquids}
\begin{tabular}{lcccc}
\toprule
\textbf{Liquid} & \textbf{$M$} & \textbf{$\rho$} & \textbf{$n$} & \textbf{$R_m$} \\
 & \si{\gram\per\mole} & \si{\gram\per\cubic\centi\meter} & RI & \si{\cubic\centi\meter\per\mole} \\
\midrule
Water      & & & & \\
Ethanol    & & & & \\
Glycerol   & & & & \\
Cyclohexane & & & & \\
\bottomrule
\end{tabular}
\end{table}

\section*{Calculations}

\subsection*{For Water}

Given:
\[
n = 1.334,\qquad M = 18~\si{\gram\per\mole},
\qquad \rho = 0.998~\si{\gram\per\cubic\centi\meter}
\]

Using the Lorentz--Lorenz equation,

\[
R_m =
\left(\frac{n^2-1}{n^2+2}\right)
\frac{M}{\rho}
\]

\[
R_m =
\left(
\frac{(1.334)^2-1}{(1.334)^2+2}
\right)
\times
\frac{18}{0.998}
\]

\[
R_m \approx 3.720~\si{\cubic\centi\meter\per\mole}
\]

\section*{Result}

The refractive indices of the given liquids are determined using an Abbe refractometer, and their molecular refractivities are calculated using the Lorentz--Lorenz equation.

\section*{Precautions}

\begin{enumerate}
    \item The prism surfaces must be completely clean and dry before every measurement.
    \item Only a few drops of sample should be used; the entire prism surface should be uniformly covered.
    \item Avoid air bubbles between the prisms.
    \item The temperature should be kept constant because refractive index is temperature-dependent.
    \item The light--dark boundary should be sharply focused before taking the reading.
    \item The cross-wire should be positioned exactly at the boundary.
    \item Volatile liquids such as acetone and methanol should be measured promptly after placing them on the prism.
    \item Do not scratch or damage the prism surfaces while cleaning them.
\end{enumerate}

\end{document}
